% Eukolia PDF subsystem test fixture.
%
% Regenerate with:
%   pdflatex -interaction=nonstopmode -halt-on-error -synctex=1 fixture.tex
% (run twice, so the table of contents and page count are stable)
%
% hyperref is not cosmetic here: it is what makes pdflatex emit a PDF outline
% (/Outlines bookmarks) and real link annotations, so the native engine's
% outline() and links() paths are exercised against genuine PDF structures
% rather than against a hand-built stub.
\documentclass[11pt]{article}
\usepackage{amsmath,amssymb}
\usepackage[margin=1in]{geometry}
\usepackage[hidelinks]{hyperref}
\title{Eukolia Native PDF Engine Fixture}
\author{Eukolia Test Suite}
\date{2026}
\begin{document}
\maketitle
\tableofcontents
\newpage
\section{Alpha Section}
\label{sec:alpha}
This paragraph exists so that the native engine has real body text to extract,
search and select. It mentions the marker AARDVARK exactly once, and it also
mentions the standalone word marmot twice, so that a case-insensitive search and
a whole-word search both have something real to find.

It carries two deliberately distinctive tokens. The first is ZEBRAQUAGGA, and it
is the only text on this page containing those five letters, so a whole-word
search for the abbreviation must fail while a substring search must succeed. The
second is QuokkaToken, whose mixed case is what the case-sensitivity assertions
are built on. Here is marmot once more.

The engine must report a sensible MediaBox for this page and must be able to
rasterise it at several scales, so the fixture deliberately contains enough
glyphs and rules that a rendered tile is never a single flat colour.

\begin{equation}
  \label{eq:euler}
  e^{i\pi} + 1 = 0
\end{equation}

A display equation above gives the renderer curves and italic glyphs to deal
with. Inline mathematics such as $\alpha^2 + \beta^2 = \gamma^2$ and a fraction
$\frac{a+b}{c+d}$ add more glyph variety within a single line.

\begin{itemize}
  \item First bullet item, deliberately long enough to exercise line wrapping.
  \item Second bullet item, mentioning the hyphenated
        compound four-and-twenty so that search can be tested across a break.
  \item Third bullet item, mentioning nothing in particular.
\end{itemize}

\subsection{Alpha Subsection}
A subsection adds another outline entry, which lets the test assert that the
outline tree nests rather than coming back flat. It also cross-references
Section~\ref{sec:beta} on page~\pageref{sec:beta}, which becomes an internal
PDF link annotation, and links to \href{https://example.com/eukolia}{an external
page}, which becomes an external link annotation.

\newpage
\section{Beta Section}
\label{sec:beta}
The second page exists so that page counting, per-page text extraction and
page-scoped search all have something real to work with. It carries a table.

\begin{tabular}{|l|r|}
  \hline
  Item & Count \\
  \hline
  Alpha & 1 \\
  Beta  & 2 \\
  Gamma & 3 \\
  \hline
\end{tabular}

\section{Gamma Section}
The third section starts on the same page as the second one, so the outline has
two entries that map onto the same PDF page.

Marker line: ZEBRAFISH.

\end{document}
